%=======================================================================
\chapter*{Conclusion Générale}
\addcontentsline{toc}{chapter}{Conclusion Générale}

\section*{Bilan du projet}
Ce projet avait pour objectif de développer un système de classification automatique des
tumeurs cérébrales à partir d'images IRM, en s'appuyant sur une démarche rigoureuse
d'ingénierie des caractéristiques et d'évaluation expérimentale.

L'approche proposée a couvert l'ensemble de la chaîne méthodologique : prétraitement des
images, extraction de descripteurs complémentaires, optimisation des paramètres, construction
d'ensembles de caractéristiques, benchmark de plusieurs algorithmes et validation statistique
des résultats.

Le travail réalisé montre qu'une approche de Machine Learning classique, lorsqu'elle est
conçue de manière structurée et systématique, peut atteindre un niveau de performance élevé
tout en conservant une bonne interprétabilité.

\section*{Résultats obtenus}
Les expérimentations menées sur les 8 ensembles de caractéristiques et les 8 modèles évalués
montrent une hiérarchie claire des performances.

La meilleure configuration est obtenue avec \textbf{SVM + Full(opt)}, qui atteint
une Accuracy de \bestAcc{} et un F1-Macro de \bestFone{} sur l'ensemble de test.
Cette performance est confirmée par les analyses de robustesse (Bootstrap, Kappa, tests
statistiques) présentées au Chapitre~7.

L'étude par branche met en évidence la forte contribution des descripteurs structurels (HOG),
tandis que la fusion multi-caractéristiques apporte un gain supplémentaire en exploitant la
complémentarité entre informations statistiques, texturales, fréquentielles et géométriques.

Au-delà du score final, le projet fournit donc une compréhension fine de l'impact de chaque
famille de caractéristiques sur la qualité de la classification.

\section*{Contributions}
Les principales contributions de ce travail peuvent être résumées comme suit :
\begin{itemize}[leftmargin=1.4em,itemsep=2pt]
  \item mise en place d'un pipeline complet et reproductible pour la classification de
  tumeurs cérébrales par IRM ;
  \item optimisation systématique des descripteurs (GLCM, LBP, DWT, HOG) via une phase dédiée
  de recherche d'hyperparamètres ;
  \item comparaison exhaustive de huit modèles de Machine Learning sur huit ensembles de
  caractéristiques, avec un leaderboard global ;
  \item validation statistique avancée des résultats (Cohen's Kappa, intervalles Bootstrap,
  McNemar, Friedman) ;
  \item analyse d'interprétabilité par importance de branches, permettant de relier les
  performances aux types d'information visuelle exploités.
\end{itemize}

Ces apports confèrent au projet une valeur à la fois applicative et méthodologique.

\section*{Perspectives}
Les perspectives identifiées s'inscrivent dans la continuité directe des résultats obtenus.

Une première direction consiste à étendre le benchmark vers des approches de
Deep Learning (VGG, ResNet, EfficientNet, Vision Transformers) afin de comparer de manière
contrôlée les représentations manuelles et profondes.

Une deuxième direction porte sur l'intégration de modèles Vision-Language capables de
produire, en plus d'une classe prédite, des explications médicales textuelles structurées.
Les stratégies de fine-tuning efficient (LoRA) et d'apprentissage guidé (Teacher Forcing)
constituent ici des leviers prometteurs.

Enfin, la perspective la plus ambitieuse est la construction d'un \textbf{framework hybride}
combinant Handcrafted Features, Deep Features et VLMs dans une architecture multimodale
unifiée. Une telle approche permettrait d'évoluer d'un classificateur performant vers un
système d'aide au diagnostic plus complet, capable de \textbf{prédire, justifier et
expliquer} ses décisions de façon exploitable par les professionnels de santé.
